\documentclass[12pt,a4paper]{article}

% ==============================================================================
% GÓI THƯ VIỆN CƠ BẢN VÀ HỖ TRỢ TIẾNG VIỆT
% ==============================================================================
\usepackage[utf8]{inputenc}
\usepackage[T5]{fontenc}
\usepackage[vietnamese]{babel}

\usepackage{amsmath, amssymb, amsthm, mathtools}
\usepackage{graphicx}
\usepackage{xcolor}
\usepackage{booktabs, multirow, makecell, tabularx, longtable, array}
\usepackage{geometry}
\geometry{
    a4paper,
    left=25mm,
    right=20mm,
    top=25mm,
    bottom=25mm
}
\usepackage{fancyhdr}
\usepackage{titlesec}
\usepackage{enumitem}
\usepackage{listings}
\usepackage{tcolorbox}
\tcbuselibrary{breakable}
\usepackage{hyperref}
\usepackage{tikz}
\usepackage{float}

% Cấu hình đường dẫn thư mục ảnh
\graphicspath{{./}{figures/}}

% Cấu hình Hyperlink
\hypersetup{
    colorlinks=true,
    linkcolor=blue!75!black,
    citecolor=red!75!black,
    urlcolor=teal!80!black
}

% Cấu hình Header & Footer
\setlength{\headheight}{15pt}
\addtolength{\topmargin}{-3pt}
\pagestyle{fancy}
\fancyhf{}
\fancyhead[L]{\small\textit{Học viện Công nghệ Bưu chính Viễn thông -- PTIT}}
\fancyhead[R]{\small\textit{Báo cáo Kỹ thuật Assignment 06 -- LTHTTM}}
\fancyfoot[C]{\thepage}
\renewcommand{\headrulewidth}{0.6pt}
\renewcommand{\footrulewidth}{0.4pt}

% Cấu hình Code Listing
\lstset{
    language=Python,
    basicstyle=\ttfamily\footnotesize,
    keywordstyle=\color{blue!80!black}\bfseries,
    commentstyle=\color{green!50!black}\itshape,
    stringstyle=\color{red!70!black},
    numbers=left,
    numberstyle=\tiny\color{gray},
    stepnumber=1,
    numbersep=8pt,
    backgroundcolor=\color{gray!5},
    frame=single,
    rulecolor=\color{gray!30},
    breaklines=true,
    breakatwhitespace=true,
    showstringspaces=false,
    tabsize=4
}

% Định dạng tiêu đề các mục
\titleformat{\section}{\large\bfseries\color{blue!85!black}}{\thesection}{1em}{}[\titlerule]
\titleformat{\subsection}{\normalsize\bfseries\color{cyan!75!black}}{\thesubsection}{1em}{}
\titleformat{\subsubsection}{\small\bfseries\color{black}}{\thesubsubsection}{1em}{}

% Khối ghi chú nổi bật
\newtcolorbox{mybox}[1]{
    colback=blue!5!white,
    colframe=blue!75!black,
    fonttitle=\bfseries,
    title=#1,
    arc=2mm,
    boxrule=1pt
}

% Khối chú thích & phân tích biểu đồ chuyên sâu
\newtcolorbox{chartbox}[1][]{
    colback=blue!2!white,
    colframe=blue!70!black,
    fonttitle=\bfseries\small,
    coltitle=white,
    arc=1.5mm,
    boxrule=0.8pt,
    left=3mm,
    right=3mm,
    top=2mm,
    bottom=2mm,
    breakable,
    title=#1
}

% Cấu hình cột bảng tự co giãn và căn lề
\newcolumntype{L}[1]{>{\raggedright\arraybackslash}p{#1}}
\newcolumntype{C}[1]{>{\centering\arraybackslash}p{#1}}
\newcolumntype{Y}{>{\centering\arraybackslash}X}

% ==============================================================================
% BẮT ĐẦU TÀI LIỆU
% ==============================================================================
\begin{document}

% ==============================================================================
% TRANG BÌA CHUẨN KHUNG VIỀN CỔ ĐIỂN VÀ HỌA TIẾT 4 GÓC THEO FILE DOCX
% ==============================================================================
\begin{titlepage}
    \begin{tikzpicture}[remember picture, overlay]
        % 1. Đường viền khung ngoài dày
        \draw[line width=2.2pt, color=black]
            ([xshift=15mm, yshift=-15mm]current page.north west)
            rectangle
            ([xshift=-15mm, yshift=15mm]current page.south east);
        
        % 2. Đường viền khung trong mảnh
        \draw[line width=0.8pt, color=black]
            ([xshift=17.5mm, yshift=-17.5mm]current page.north west)
            rectangle
            ([xshift=-17.5mm, yshift=17.5mm]current page.south east);
    \end{tikzpicture}

    \centering
    \vspace*{0.2cm}
    {\large \textbf{HỌC VIỆN CÔNG NGHỆ BƯU CHÍNH VIỄN THÔNG}}\\[0.25cm]
    {\large \textbf{KHOA CÔNG NGHỆ THÔNG TIN 1}}\\[0.2cm]
    \rule{5cm}{1pt}\\[0.9cm]

    % Logo PTIT
    \begin{center}
        \IfFileExists{figures/logo_ptit.png}{%
            \includegraphics[width=3.4cm, height=4.2cm, keepaspectratio]{figures/logo_ptit.png}
        }{%
            \begin{tcolorbox}[colback=red!5!white, colframe=red!75!black, width=3.8cm, height=3.8cm, arc=2mm, halign=center, valign=center]
                \Large \textbf{\color{red!80!black}PTIT}\\[0.1cm]
                \footnotesize \textbf{\color{black}LOGO}
            \end{tcolorbox}
        }
    \end{center}

    \vspace{0.6cm}

    {\Large \textbf{TIỂU LUẬN CUỐI KỲ}}\\[0.3cm]
    {\large \textbf{MÔN HỌC: LẬP TRÌNH HỆ THỐNG THÔNG MINH}}\\[0.5cm]

    \vspace{0.8cm}

    \begin{flushleft}
        \hspace{1.6cm}
        \begin{tabular}{p{5.0cm} p{0.5cm} l}
            \textbf{Giảng viên hướng dẫn} & : & \textbf{PGS.TS. Trần Đình Quế} \\[0.3cm]
            \textbf{Lớp}                  & : & \textbf{D23CTPM01} \\[0.3cm]
            \textbf{Nhóm học phần}        & : & \textbf{CT01} \\[0.3cm]
            \textbf{Thành viên thực hiện} & : & \textbf{Nguyễn Nam Hải} \\[0.3cm]
            \textbf{Mã sinh viên}         & : & \textbf{B23DCCN277} \\[0.3cm]
        \end{tabular}
    \end{flushleft}

    \vfill
    {\large \textbf{Hà Nội -- 2026}}
    \vspace*{0.3cm}
\end{titlepage}

% ==============================================================================
% MỤC LỤC BÁO CÁO KỸ THUẬT
% ==============================================================================
\newpage
\renewcommand{\contentsname}{MỤC LỤC TIỂU LUẬN CUỐI KỲ}
\tableofcontents
\newpage

% ==============================================================================
% DANH MỤC HÌNH VẼ & BẢNG BIỂU
% ==============================================================================
\phantomsection
\addcontentsline{toc}{section}{DANH MỤC HÌNH VẼ}
\listoffigures

\vspace{0.8cm}

\phantomsection
\addcontentsline{toc}{section}{DANH MỤC BẢNG BIỂU}
\listoftables

\newpage

% ==============================================================================
% TÓM TẮT BÁO CÁO (ABSTRACT)
% ==============================================================================
\phantomsection
\addcontentsline{toc}{section}{TÓM TẮT BÁO CÁO (ABSTRACT)}
\section*{TÓM TẮT BÁO CÁO (ABSTRACT)}
\addcontentsline{toc}{section}{TÓM TẮT BÁO CÁO (ABSTRACT)}

Tiểu luận môn học Lập trình hệ thống thông minh trình bày tổng quan về Trí tuệ nhân tạo, từ lịch sử phát triển cho đến các kỹ thuật Học máy (Machine Learning) cơ bản, Mạng nơ-ron Tích chập (CNN) và Mạng nơ-ron Hồi quy (RNN). Báo cáo thực hiện khảo sát các tập dữ liệu, xây dựng mô hình từ đầu (from scratch) cho đến sử dụng các framework hiện đại như Keras và PyTorch, đồng thời so sánh đánh giá hiệu năng và triển khai ứng dụng thực tế.

\newpage

\section{Mở đầu}
\subsection{Lý do chọn đề tài}
Trong kỷ nguyên công nghệ số hiện đại, Trí tuệ Nhân tạo (Artificial Intelligence - AI) đã và đang tạo ra những bước ngoặt mang tính cách mạng trong hầu hết mọi lĩnh vực của đời sống con người. Từ những hệ thống gợi ý sản phẩm thông minh trên các nền tảng thương mại điện tử, hệ thống chẩn đoán y khoa tự động, xe tự lái, cho đến các mô hình ngôn ngữ lớn (Large Language Models) như ChatGPT, AI không còn là một khái niệm viễn tưởng mà đã trở thành một công cụ thiết yếu. Sự phát triển bùng nổ của lượng dữ liệu khổng lồ (Big Data) kết hợp với năng lực tính toán ngày càng vượt trội của các phần cứng đồ họa (GPU/TPU) đã mở ra một kỷ nguyên vàng cho các thuật toán Học máy (Machine Learning - ML) và Học sâu (Deep Learning - DL).

Đằng sau những ứng dụng hào nhoáng đó là một nền tảng toán học và khoa học máy tính vô cùng đồ sộ. Để có thể làm chủ và ứng dụng AI một cách hiệu quả, việc hiểu rõ bản chất cốt lõi của các thuật toán, từ những phương pháp học máy thống kê cơ bản cho đến các kiến trúc mạng nơ-ron phức tạp như Mạng nơ-ron Tích chập (Convolutional Neural Networks - CNN) hay Mạng nơ-ron Hồi quy (Recurrent Neural Networks - RNN), là một yêu cầu bắt buộc đối với bất kỳ kỹ sư công nghệ thông tin nào. Quá trình học tập không chỉ dừng lại ở việc gọi các hàm có sẵn từ các thư viện cấp cao (API-driven development) mà còn phải đi sâu vào việc triển khai từ đầu (from scratch) để thấu hiểu cách thức tối ưu hóa hàm mất mát (loss function), sự lan truyền ngược của đạo hàm (backpropagation), và cách thức dữ liệu được biến đổi qua các không gian vector đa chiều.

\subsection{Mục tiêu nghiên cứu}
Nghiên cứu này được thực hiện với các mục tiêu cốt lõi như sau:
\begin{itemize}
    \item \textbf{Hệ thống hóa cơ sở lý thuyết:} Trình bày một cách logic và khoa học quá trình tiến hóa của các hệ thống AI, từ những hệ chuyên gia sơ khai cho đến sự bùng nổ của mạng nơ-ron nhân tạo. Phân tích sâu các khái niệm toán học làm nền tảng cho Học máy.
    \item \textbf{Thực nghiệm hóa thuật toán cơ bản:} Xây dựng và đánh giá các mô hình Phân loại và Hồi quy sử dụng Multi-Layer Perceptron (MLP) trên các tập dữ liệu thực tế thuộc các miền lĩnh vực đa dạng như y tế, bất động sản và thương mại điện tử.
    \item \textbf{Khai phá kiến trúc Mạng Tích chập (CNN):} Phân tích cơ chế trích xuất đặc trưng không gian thông qua phép tích chập (convolution) và gộp (pooling). Đánh giá tính hiệu quả trên ảnh màu, ảnh xám và cả dữ liệu bảng 1D.
    \item \textbf{Giải quyết bài toán Chuỗi thời gian (RNN):} Phân tích cấu trúc hồi quy, thuật toán BPTT và ứng dụng dự báo chuỗi thời gian liên tục trên dữ liệu tài chính vĩ mô.
    \item \textbf{Đánh giá đối chuẩn (Benchmarking):} Mọi thuật toán đều được triển khai song song trên 3 nền tảng: NumPy (From scratch), Keras (High-level API) và PyTorch (Dynamic Computation Graph) để cung cấp cái nhìn khách quan về hiệu năng và độ phức tạp mã nguồn.
\end{itemize}

\subsection{Phạm vi và đối tượng nghiên cứu}
\textbf{Đối tượng nghiên cứu:}
Nghiên cứu tập trung vào 3 cấu trúc mạng nơ-ron tiêu biểu: MLP, CNN, và RNN. Song song đó là các thuật toán tối ưu hóa như Gradient Descent (BGD, SGD, Mini-batch) và các phương pháp tinh chỉnh siêu tham số, chuẩn hóa dữ liệu.

\textbf{Phạm vi thực nghiệm (Datasets):}
Để đảm bảo tính khách quan, tiểu luận sử dụng 9 tập dữ liệu khác biệt về bản chất:
\begin{itemize}
    \item Dữ liệu dạng bảng: Bệnh tiểu đường Pima (Diabetes), Giá nhà California (Housing), Đánh giá thời trang (E-commerce).
    \item Dữ liệu thị giác: Ảnh vật thể CIFAR-10, Ảnh chữ số viết tay MNIST, và một biến thể CNN 1D cho dữ liệu bảng.
    \item Dữ liệu tuần tự (Time-series): Giá Vàng (Gold Price), Cổ phiếu Amazon (AMZN), và dữ liệu khí tượng Jena Climate.
\end{itemize}

\subsection{Cấu trúc của Tiểu luận}
Báo cáo được cấu trúc thành 4 chương chính. Chương 1 tóm tắt lịch sử AI từ thập niên 1950 đến nay. Chương 2 đi sâu vào các mô hình Machine Learning cơ bản (Perceptron đa tầng). Chương 3 tập trung phân tích Mạng nơ-ron Tích chập (CNN) và khả năng trích xuất đặc trưng không gian. Cuối cùng, Chương 4 trình bày về Mạng nơ-ron Hồi quy (RNN) trong việc giải quyết dữ liệu chuỗi thời gian.

\section{Lịch sử phát triển của Trí tuệ nhân tạo}

Trí tuệ nhân tạo (Artificial Intelligence - AI) là một trong những lĩnh vực nghiên cứu năng động và phát triển nhanh nhất trong khoa học máy tính. Mặc dù thuật ngữ "Trí tuệ nhân tạo" mới chỉ chính thức được đưa ra vào giữa thế kỷ 20, nhưng ước mơ về việc tạo ra những cỗ máy có khả năng suy nghĩ và hành động như con người đã tồn tại từ thời cổ đại qua những câu chuyện thần thoại, truyền thuyết và các tác phẩm cơ khí sơ khai. Sự phát triển của AI không phải là một đường thẳng đi lên, mà là một chuỗi các giai đoạn với những bước tiến đột phá đan xen với những kỳ vọng bị phá vỡ, thường được gọi là các "Mùa đông AI". Chương này sẽ đi sâu vào lịch sử hình thành, sự phát triển qua các thời kỳ, những cột mốc quan trọng và trạng thái hiện tại của Trí tuệ nhân tạo, đồng thời xem xét chi tiết về các cơ sở toán học và thuật toán trong những thời kỳ đầu.

Sự hình thành của AI bắt nguồn từ các lĩnh vực triết học, toán học, tâm lý học, thần kinh học và khoa học máy tính. Những câu hỏi cốt lõi như "Tâm trí hoạt động như thế nào?", "Làm thế nào một hệ thống vật lý có thể tạo ra trí thông minh?" đã thôi thúc các nhà nghiên cứu tìm kiếm cách mô phỏng các quá trình nhận thức của con người bằng các thiết bị tính toán cơ học và điện tử. Trong thế kỷ 20, sự ra đời của máy tính điện tử đã cung cấp nền tảng phần cứng cần thiết để biến những lý thuyết về trí tuệ nhân tạo thành hiện thực. 

Lịch sử AI có thể được chia thành nhiều giai đoạn chính: sự khởi đầu khái niệm với Alan Turing, Hội nghị Dartmouth huyền thoại đánh dấu sự ra đời của lĩnh vực học thuật, những thập kỷ của sự lạc quan tột độ (và sự sụp đổ sau đó), sự trỗi dậy của các hệ chuyên gia, sự trở lại của mạng nơ-ron nhân tạo nhờ thuật toán lan truyền ngược, và cuối cùng là kỷ nguyên bùng nổ của Học sâu (Deep Learning) và các mô hình ngôn ngữ lớn (Large Language Models) mà chúng ta đang chứng kiến ngày nay.

\newpage

\subsection{Sự ra đời của Trí tuệ nhân tạo (Trước 1956 và Hội nghị Dartmouth)}

Thời kỳ trước năm 1956 đánh dấu những bước chuẩn bị quan trọng về mặt lý thuyết và khái niệm cho sự ra đời của AI. Một trong những cột mốc quan trọng nhất là bài báo kinh điển "Computing Machinery and Intelligence" được xuất bản năm 1950 bởi nhà toán học, logic học và mật mã học vĩ đại người Anh, Alan Turing. Trong bài báo này, Turing đã đặt ra một câu hỏi nền tảng: "Máy móc có thể suy nghĩ không?" (Can machines think?). Để tránh những tranh cãi triết học phức tạp về định nghĩa của "suy nghĩ", Turing đã đề xuất một bài kiểm tra thực nghiệm, sau này được biết đến rộng rãi với tên gọi "Phép thử Turing" (Turing Test).

Phép thử Turing được thiết kế như một "Trò chơi bắt chước" (Imitation Game), trong đó một người thẩm vấn (interrogator) sẽ giao tiếp với một con người và một máy tính thông qua văn bản (text) ở hai phòng tách biệt. Nếu người thẩm vấn không thể phân biệt được đâu là máy tính và đâu là con người sau một khoảng thời gian giao tiếp, máy tính đó được coi là đã vượt qua phép thử và có khả năng thể hiện trí thông minh tương đương con người. Bài báo của Turing không chỉ giới thiệu Phép thử Turing mà còn thảo luận về nhiều vấn đề then chốt của AI như học máy (machine learning), thuật toán di truyền (genetic algorithms) và học củng cố (reinforcement learning).

Bên cạnh Alan Turing, những nghiên cứu đầu tiên về mạng nơ-ron nhân tạo cũng được tiến hành trong giai đoạn này. Năm 1943, Warren McCulloch (nhà thần kinh học) và Walter Pitts (nhà logic học) đã xuất bản bài báo "A Logical Calculus of the Ideas Immanent in Nervous Activity", đề xuất một mô hình toán học đơn giản đầu tiên của nơ-ron sinh học. Họ chứng minh rằng các mạng bao gồm các nơ-ron nhân tạo này có thể tính toán bất kỳ hàm logic nào. 

Về mặt toán học, mô hình McCulloch-Pitts có thể được biểu diễn như sau. Mỗi nơ-ron nhận các tín hiệu đầu vào nhị phân $x_i \in \{0, 1\}$. Trọng số của các kết nối, $w_i$, có thể là kích thích (excitatory, $w_i > 0$) hoặc ức chế (inhibitory, $w_i < 0$). Trạng thái đầu ra $y$ của nơ-ron được xác định bởi một hàm bước (step function):

\begin{equation}
y = \begin{cases} 
1 & \text{nếu } \sum_{i=1}^{n} w_i x_i \geq \theta \\
0 & \text{nếu } \sum_{i=1}^{n} w_i x_i < \theta 
\end{cases}
\end{equation}

trong đó $\theta$ là ngưỡng kích hoạt (threshold). Mặc dù đơn giản, mô hình này đã đặt nền móng lý thuyết quan trọng cho cấu trúc của các mạng nơ-ron hiện đại.

Năm 1956 được coi là năm sinh chính thức của lĩnh vực Trí tuệ nhân tạo với "Dự án Nghiên cứu Mùa hè về Trí tuệ nhân tạo tại Dartmouth" (Dartmouth Summer Research Project on Artificial Intelligence). Hội nghị này được tổ chức bởi John McCarthy, Marvin Minsky, Nathaniel Rochester và Claude Shannon. John McCarthy đã chính thức đề xuất việc sử dụng thuật ngữ "Artificial Intelligence" (Trí tuệ nhân tạo) trong đề xuất của dự án. Đề xuất của hội nghị có tuyên bố đầy tham vọng: "Nghiên cứu sẽ được tiến hành dựa trên phỏng đoán rằng mọi khía cạnh của việc học tập hoặc bất kỳ đặc điểm nào khác của trí thông minh đều có thể được mô tả chính xác đến mức một cỗ máy có thể được chế tạo để mô phỏng nó". 

Hội nghị Dartmouth quy tụ những bộ óc sáng giá nhất thời bấy giờ, trong đó có Allen Newell và Herbert A. Simon, những người đã trình bày chương trình "Logic Theorist" - chương trình máy tính đầu tiên được thiết kế để bắt chước kỹ năng giải quyết vấn đề của con người, và đã chứng minh thành công nhiều định lý trong cuốn Principia Mathematica. Sự kiện này đã thống nhất một cộng đồng nhỏ các nhà nghiên cứu xung quanh một mục tiêu chung và khởi xướng nhiều thập kỷ nghiên cứu chuyên sâu về AI.

\newpage

\subsection{Kỳ vọng ban đầu và Mùa đông AI đầu tiên (1956 - 1980)}

Những năm tiếp theo sau Hội nghị Dartmouth chứng kiến một sự lạc quan tột độ và những khoản đầu tư khổng lồ từ các cơ quan chính phủ, đặc biệt là DARPA (Defense Advanced Research Projects Agency) của Hoa Kỳ. Những người tiên phong trong lĩnh vực AI tin rằng máy móc có trí thông minh ngang ngửa con người sẽ xuất hiện trong vòng một vài thập kỷ.

Trong giai đoạn này, nhiều chương trình máy tính đáng kinh ngạc đã được phát triển. Ví dụ tiêu biểu là General Problem Solver (GPS) do Newell và Simon phát triển năm 1957. GPS được thiết kế để giải quyết các vấn đề đa dạng theo cách tư duy của con người, sử dụng phương pháp phân tích phương tiện-mục đích (means-ends analysis) để giảm thiểu sự khác biệt giữa trạng thái hiện tại và trạng thái đích. Mặc dù GPS giải quyết tốt các bài toán logic đơn giản, nó gặp khó khăn khi đối mặt với các vấn đề trong thế giới thực đòi hỏi kiến thức nền rộng lớn.

Một thành tựu nổi bật khác là chương trình ELIZA, do Joseph Weizenbaum tại MIT phát triển từ năm 1964 đến 1966. ELIZA là một trong những chương trình xử lý ngôn ngữ tự nhiên đầu tiên, được thiết kế để mô phỏng một nhà trị liệu tâm lý Rogerian. Mặc dù ELIZA chỉ sử dụng các kỹ thuật đối sánh mẫu (pattern matching) và thay thế văn bản đơn giản mà không thực sự "hiểu" ngữ nghĩa, nó đã tạo ra ảo giác mạnh mẽ về trí thông minh và sự đồng cảm, khiến nhiều người dùng tin rằng họ đang trò chuyện với một thực thể có trí óc.

Tuy nhiên, sự lạc quan quá mức đã dẫn đến những hứa hẹn phi thực tế. Khi những khó khăn thực sự của việc giải quyết các bài toán AI phức tạp lộ diện, tiến độ chậm lại đáng kể. Các hệ thống AI thời kỳ này chủ yếu dựa trên logic biểu tượng (symbolic logic) và các quy tắc tìm kiếm heuristic. Mặc dù phương pháp này có hiệu quả với các "thế giới khối" (micro-worlds) bị giới hạn, nó hoàn toàn thất bại khi đối mặt với "sự bùng nổ tổ hợp" (combinatorial explosion) trong các vấn đề thế giới thực quy mô lớn. Số lượng các trạng thái có thể xảy ra trong các bài toán thực tế tăng lên theo cấp số nhân, vượt xa khả năng tính toán của máy tính và bộ nhớ phần cứng lúc bấy giờ.

Hơn nữa, một rào cản lớn khác là "nghịch lý Moravec" (Moravec's paradox). Nghịch lý này phát biểu rằng các nhiệm vụ tư duy cấp cao (như giải toán logic, chơi cờ) đòi hỏi rất ít tính toán, trong khi các nhiệm vụ nhận thức và vận động cấp thấp (như nhận dạng hình ảnh, di chuyển trong không gian, hiểu ngôn ngữ giao tiếp hàng ngày) lại đòi hỏi những nguồn lực tính toán khổng lồ. 

Báo cáo Lighthill (Lighthill Report) năm 1973 tại Anh đã đưa ra một đánh giá rất bi quan về triển vọng của AI, chỉ trích lĩnh vực này vì không đạt được những mục tiêu vĩ đại đã hứa hẹn. Báo cáo này dẫn đến việc chính phủ Anh cắt giảm nghiêm trọng tài trợ cho các nghiên cứu AI tại hầu hết các trường đại học. Cùng thời điểm đó ở Mỹ, DARPA cũng rút bớt tài trợ cho nghiên cứu AI cơ bản do thất vọng về kết quả của dự án dịch máy tự động. Sự cắt giảm tài trợ hàng loạt và sự suy giảm niềm tin này đã mở ra "Mùa đông AI đầu tiên" (The First AI Winter), một giai đoạn kéo dài từ giữa những năm 1970 đến khoảng năm 1980, đặc trưng bởi sự đình trệ và thiếu nguồn lực trong nghiên cứu trí tuệ nhân tạo.

\newpage

\subsection{Hệ chuyên gia và Mùa đông AI thứ hai (1980 - 1993)}

Vào đầu những năm 1980, AI bắt đầu hồi sinh nhờ một hướng đi mới và thực dụng hơn: các Hệ chuyên gia (Expert Systems). Thay vì cố gắng tạo ra một trí thông minh tổng quát (Artificial General Intelligence - AGI) như trước đây, các nhà nghiên cứu nhận ra rằng họ có thể tạo ra các hệ thống hữu ích bằng cách giới hạn phạm vi ứng dụng vào các miền tri thức hẹp và cụ thể.

Một hệ chuyên gia là một hệ thống máy tính mô phỏng khả năng ra quyết định của một chuyên gia con người trong một lĩnh vực cụ thể. Hệ thống bao gồm hai thành phần chính:
\begin{itemize}
    \item \textbf{Cơ sở tri thức (Knowledge Base):} Chứa đựng kiến thức chuyên môn dưới dạng các tập luật (thường là các câu điều kiện IF-THEN) và sự kiện do các chuyên gia cung cấp.
    \item \textbf{Động cơ suy diễn (Inference Engine):} Các thuật toán logic áp dụng các luật trong cơ sở tri thức vào dữ liệu đầu vào để suy luận và đưa ra câu trả lời hoặc khuyến nghị.
\end{itemize}

Hệ chuyên gia đầu tiên thành công vang dội là R1 (sau này được đổi tên thành XCON), được phát triển bởi John McDermott tại Đại học Carnegie Mellon vào năm 1980 cho Tập đoàn Thiết bị Kỹ thuật số (Digital Equipment Corporation - DEC). XCON tự động hóa quá trình cấu hình các hệ thống máy tính VAX của DEC dựa trên yêu cầu của khách hàng, giúp DEC tiết kiệm hàng chục triệu USD mỗi năm bằng cách giảm thiểu lỗi trong các đơn đặt hàng. Một ví dụ nổi tiếng khác là MYCIN, phát triển tại Đại học Stanford, được thiết kế để chẩn đoán các bệnh nhiễm khuẩn máu và đề xuất liệu pháp kháng sinh. MYCIN có độ chính xác chẩn đoán vượt trội so với nhiều bác sĩ nội trú.

Sự thành công thương mại của các hệ chuyên gia đã kích hoạt một làn sóng đầu tư mới khổng lồ vào AI từ các tập đoàn công nghiệp. Một ngành công nghiệp hoàn toàn mới xoay quanh phần cứng và phần mềm cho AI đã ra đời, trong đó đáng chú ý nhất là thị trường máy tính Lisp (Lisp machines) – những máy tính được thiết kế đặc biệt để tối ưu hóa việc chạy ngôn ngữ lập trình Lisp phổ biến trong AI thời đó, do các công ty như Symbolics và Lisp Machines Inc. sản xuất. Chính phủ Nhật Bản cũng công bố "Dự án Máy tính Thế hệ thứ năm" (Fifth Generation Computer Systems) với khoản ngân sách hàng trăm triệu đô la, nhằm phát triển siêu máy tính có khả năng suy luận logic khổng lồ.

Tuy nhiên, thời kỳ hoàng kim của hệ chuyên gia không kéo dài. Đến cuối những năm 1980, các công ty bắt đầu nhận ra những giới hạn chí mạng của hướng tiếp cận này:
\begin{enumerate}
    \item \textbf{Khó khăn trong việc thu nhận tri thức (Knowledge Acquisition Bottleneck):} Việc trích xuất kiến thức từ các chuyên gia con người và mã hóa chúng thành các tập luật là quá trình tốn cực kỳ nhiều thời gian, công sức và tốn kém.
    \item \textbf{Tính giòn gãy (Brittleness):} Hệ chuyên gia không có "ý thức" thông thường. Khi đối mặt với một tình huống nằm ngoài các quy tắc đã được lập trình sẵn, hệ thống sẽ thất bại hoàn toàn thay vì suy thoái dần (graceful degradation) như trí tuệ con người.
    \item \textbf{Chi phí bảo trì cao:} Khi lĩnh vực chuyên môn thay đổi, cơ sở tri thức phải được cập nhật thủ công, điều này rất phức tạp đối với các hệ thống có hàng chục ngàn tập luật.
    \item \textbf{Sự sụp đổ của phần cứng chuyên dụng:} Sự phát triển thần tốc của các máy tính cá nhân (PC) và máy trạm (workstation) thông thường từ Apple và IBM đã khiến những chiếc Lisp machines đắt đỏ trở nên lỗi thời. Các máy PC thông thường có thể chạy phần mềm Lisp nhanh hơn và rẻ hơn các máy Lisp chuyên dụng.
\end{enumerate}

Sự thất vọng của giới công nghiệp khi các dự án hệ chuyên gia thất bại trong việc đáp ứng kỳ vọng, cùng với sự kết thúc tài trợ của chính phủ, đã dẫn đến sự sụp đổ của thị trường máy tính Lisp vào năm 1987 và gây ra "Mùa đông AI thứ hai". Thời kỳ này thậm chí còn tồi tệ hơn lần đầu tiên, với cụm từ "Trí tuệ nhân tạo" gần như trở thành một từ cấm kỵ trong các hồ sơ xin tài trợ nghiên cứu hoặc trong kinh doanh, buộc các nhà nghiên cứu phải ngụy trang công việc của họ dưới các tên gọi khác như "Tin học học", "Hệ thống thông minh", hoặc "Học máy".

\newpage

\subsection{Sự trỗi dậy của Thuyết kết nối và Học máy (1990s - 2010s)}

Trong khi logic biểu tượng (Symbolic AI hay GOFAI - Good Old-Fashioned AI) thống trị phần lớn thời gian trước những năm 1990, một trường phái nghiên cứu khác, được gọi là "Thuyết kết nối" (Connectionism), vẫn âm thầm phát triển. Thuyết kết nối cố gắng mô phỏng mạng lưới sinh học của não bộ thông qua các Mạng Nơ-ron Nhân tạo (Artificial Neural Networks - ANNs).

Vào năm 1958, Frank Rosenblatt đã phát minh ra Perceptron, một mô hình nơ-ron đơn giản có khả năng học hỏi qua dữ liệu để thực hiện phân loại nhị phân. Công thức toán học cốt lõi cho đầu ra $\hat{y}$ của một Perceptron cho một mẫu dữ liệu đầu vào có đặc trưng biểu diễn dưới dạng vector $x = [x_1, x_2, \dots, x_n]^T$ được định nghĩa bằng hàm dấu (sign function):

\begin{equation}
\hat{y} = \text{sign}\left(\sum_{i=1}^{n} w_i x_i + b\right) = \text{sign}(\mathbf{w}^T \mathbf{x} + b)
\end{equation}

trong đó $\mathbf{w}$ là vector trọng số (weights) và $b$ là hệ số độ lệch (bias). Perceptron có thuật toán cập nhật trọng số vô cùng đơn giản dựa trên sai số giữa dự đoán $\hat{y}$ và nhãn thực tế $y$:

\begin{equation}
\mathbf{w}_{new} = \mathbf{w}_{old} + \eta (y - \hat{y})\mathbf{x}
\end{equation}
với $\eta$ là tốc độ học (learning rate). 

Rosenblatt rất tự tin về phát minh của mình, nhưng vào năm 1969, Marvin Minsky và Seymour Papert đã xuất bản cuốn sách "Perceptrons", trong đó họ chứng minh nghiêm ngặt bằng toán học rằng các Perceptron đơn lớp (single-layer perceptron) không thể học được các hàm phân tách phi tuyến tính cơ bản, điển hình là bài toán logic XOR (Exclusive OR). Cuốn sách này đã gây ảnh hưởng lớn đến mức khiến hướng nghiên cứu mạng nơ-ron gần như đóng băng trong suốt nhiều năm.

Phải đến giữa những năm 1980, sự quan tâm đến mạng nơ-ron mới bắt đầu hồi sinh nhờ việc phổ biến hóa thuật toán Lan truyền ngược (Backpropagation). Mặc dù thuật toán này đã được phát minh độc lập bởi nhiều nhà nghiên cứu từ những năm 1960 (như Paul Werbos), nhưng nó chỉ thực sự trở nên nổi tiếng thông qua bài báo đột phá năm 1986 của David Rumelhart, Geoffrey Hinton và Ronald Williams. Lan truyền ngược cung cấp một phương pháp hiệu quả để tính toán gradient (đạo hàm riêng) của hàm mất mát theo từng trọng số trong một mạng nơ-ron đa lớp (Multi-layer Perceptron - MLP), từ đó cho phép sử dụng thuật toán Gradient Descent để tối ưu hóa trọng số:

Giả sử hàm mất mát là $E$, theo quy tắc chuỗi (chain rule) trong giải tích, đạo hàm của $E$ theo trọng số $w_{ij}$ kết nối nơ-ron $j$ ở lớp ẩn với nơ-ron $i$ ở lớp sau là:
\begin{equation}
\frac{\partial E}{\partial w_{ij}} = \frac{\partial E}{\partial o_i} \cdot \frac{\partial o_i}{\partial net_i} \cdot \frac{\partial net_i}{\partial w_{ij}}
\end{equation}
trong đó $net_i$ là tổng đầu vào của nơ-ron $i$, và $o_i = f(net_i)$ là đầu ra sau hàm kích hoạt $f$. Phát hiện này chứng minh rằng mạng nhiều lớp hoàn toàn có khả năng học các hàm phi tuyến tính cực kỳ phức tạp (như XOR), bác bỏ lời chỉ trích của Minsky trước đây.

Năm 1989, Yann LeCun cùng các cộng sự tại AT\&T Bell Labs đã áp dụng thành công thuật toán Backpropagation vào Mạng Nơ-ron Tích chập (Convolutional Neural Networks - CNNs) để nhận dạng chữ số viết tay trên các tấm séc ngân hàng, tạo ra hệ thống LeNet.

Bất chấp những thành công này, trong thập niên 1990 và đầu những năm 2000, mạng nơ-ron lại một lần nữa bị lu mờ trước sự bùng nổ của các thuật toán Học máy (Machine Learning) truyền thống dựa trên thống kê, nổi bật là Support Vector Machines (SVM), Random Forests, và Hidden Markov Models (HMM). Các thuật toán SVM có nền tảng lý thuyết tối ưu hóa toán học vô cùng vững chắc và hoạt động cực kỳ hiệu quả trên các tập dữ liệu nhỏ và trung bình mà không cần tinh chỉnh quá nhiều siêu tham số (hyperparameters). Trong thời kỳ này, trí tuệ nhân tạo trở nên khoa học hơn, có cấu trúc tốt hơn, và dần dần được tích hợp vào các hệ thống khai phá dữ liệu (data mining) và các ứng dụng thương mại một cách thầm lặng.

\newpage

\subsection{Sự bùng nổ của Học sâu (2011 - nay)}

Bước ngoặt thực sự của kỷ nguyên trí tuệ nhân tạo hiện đại diễn ra vào đầu những năm 2010. Mặc dù các cấu trúc mạng nơ-ron đa lớp và thuật toán tối ưu hóa đã có từ nhiều thập kỷ trước, chúng thường gặp phải bài toán vanishing gradient (gradient biến mất) khiến mạng sâu không thể học được, hoặc overfitting do không có đủ dữ liệu. Ba yếu tố mang tính cách mạng đã cùng hội tụ trong giai đoạn này để tạo ra sự bùng nổ của Học sâu (Deep Learning):

\textbf{1. Sự bùng nổ của Dữ liệu lớn (Big Data):} Sự phát triển của Internet, mạng xã hội, và các thiết bị di động đã tạo ra một lượng dữ liệu khổng lồ (hình ảnh, văn bản, âm thanh). Việc có các tập dữ liệu được gán nhãn lớn là điều kiện tiên quyết cho Deep Learning. Điển hình nhất là tập dữ liệu ImageNet do nhóm của giáo sư Fei-Fei Li xây dựng, chứa hàng triệu hình ảnh được gán nhãn thủ công cho hơn một ngàn danh mục đối tượng.

\textbf{2. Sức mạnh phần cứng (GPU):} Các mạng nơ-ron sâu yêu cầu hàng tỷ phép tính ma trận dấu phẩy động. Các Bộ vi xử lý Đồ họa (GPU) do NVIDIA sản xuất, ban đầu được thiết kế để xử lý đồ họa trò chơi điện tử, với kiến trúc xử lý song song hàng nghìn lõi (CUDA cores), tỏ ra hoàn toàn lý tưởng cho việc huấn luyện mạng nơ-ron. Một mô hình từng mất hàng tháng để huấn luyện trên CPU nay có thể hoàn thành trong vài ngày trên GPU.

\textbf{3. Cải tiến Thuật toán và Hàm kích hoạt:} Sự thay thế các hàm kích hoạt truyền thống như Sigmoid hay Tanh bằng hàm ReLU (Rectified Linear Unit), $f(x) = \max(0, x)$, đã giúp giải quyết triệt để vấn đề đạo hàm biến mất. Cùng với đó là các kỹ thuật như Dropout (loại bỏ ngẫu nhiên nơ-ron để chống overfitting) và Batch Normalization.

Sự kiện làm "chấn động" cộng đồng học máy thế giới là cuộc thi phân loại hình ảnh ImageNet (ILSVRC) năm 2012. Nhóm nghiên cứu của giáo sư Geoffrey Hinton tại Đại học Toronto, bao gồm Alex Krizhevsky và Ilya Sutskever, đã giới thiệu AlexNet, một mạng CNN sâu chạy trên GPU. AlexNet đã nghiền nát tất cả các phương pháp thị giác máy tính truyền thống bằng cách giảm tỷ lệ lỗi xuống chỉ còn 15.3\% (so với mức 26.2\% của đối thủ đứng thứ hai). Cột mốc này chứng minh sức mạnh áp đảo của Học sâu trong việc tự động trích xuất các đặc trưng phân cấp từ dữ liệu thô, mở ra kỷ nguyên phục hưng cho AI.

Sau AlexNet, Học sâu tiếp tục càn quét mọi lĩnh vực. Năm 2016, hệ thống AlphaGo do DeepMind phát triển đã đánh bại Lee Sedol, một trong những kỳ thủ cờ vây vĩ đại nhất thế giới, bằng cách kết hợp Học sâu với Học củng cố (Deep Reinforcement Learning). Đây là một cột mốc lịch sử vì cờ vây được cho là trò chơi có không gian trạng thái vô tận (lớn hơn số lượng nguyên tử trong vũ trụ quan sát được) và không thể giải bằng phương pháp tìm kiếm vét cạn bằng máy tính.

\textbf{Cuộc cách mạng Transformer trong Xử lý ngôn ngữ tự nhiên:}
Năm 2017, các nhà nghiên cứu tại Google đã xuất bản bài báo định hình lại toàn bộ lĩnh vực NLP: "Attention Is All You Need". Bài báo giới thiệu kiến trúc \textbf{Transformer}, hoàn toàn loại bỏ các cấu trúc mạng tuần hoàn (RNN/LSTM) trước đó vốn nổi tiếng vì khó song song hóa và hay quên ngữ cảnh xa. Cốt lõi của kiến trúc Transformer là cơ chế \textit{Self-Attention} (tự chú ý), cho phép mô hình tính toán sự tương quan giữa các từ trong câu một cách đồng thời, bất kể vị trí của chúng.

Công thức toán học của cơ chế Scaled Dot-Product Attention được biểu diễn rất thanh lịch:

\begin{equation}
\text{Attention}(Q, K, V) = \text{softmax}\left(\frac{QK^T}{\sqrt{d_k}}\right)V
\end{equation}

trong đó $Q$ (Query), $K$ (Key), và $V$ (Value) là các ma trận đặc trưng được biến đổi từ chuỗi đầu vào, và $d_k$ là kích thước của các vector key, đóng vai trò như một hệ số tỷ lệ để tránh việc hàm softmax bị bão hòa. Cơ chế Multi-Head Attention sau đó kết hợp nhiều không gian biểu diễn khác nhau, cho phép mô hình thu thập thông tin vô cùng phong phú về ngữ pháp và ngữ nghĩa.

Dựa trên kiến trúc Transformer, OpenAI đã phát triển dòng mô hình GPT (Generative Pre-trained Transformer), bao gồm mô hình ngôn ngữ lớn (LLMs). Đồng thời, Google phát triển BERT. Phương pháp tiếp cận "Học tự giám sát" (Self-supervised learning) cho phép các mô hình này học hỏi ngữ nghĩa từ hàng tỷ tài liệu trên Internet bằng cách dự đoán các từ bị che khuất hoặc dự đoán từ tiếp theo mà không cần dữ liệu được gán nhãn của con người.

\newpage

\subsection{Tình trạng hiện tại và Tương lai của Trí tuệ nhân tạo}

Chúng ta đang sống trong kỷ nguyên của Trí tuệ nhân tạo Tạo sinh (Generative AI), với sự ra đời của các Mô hình Ngôn ngữ Lớn (LLMs) như GPT-4, Gemini (Google), Claude (Anthropic), và LLaMA (Meta). Sự xuất hiện của ChatGPT vào cuối năm 2022 được ví như "Khoảnh khắc iPhone" của AI. Lần đầu tiên, một hệ thống AI sở hữu khả năng suy luận, lập trình, viết luận văn sáng tạo, phân tích dữ liệu, và đối thoại lưu loát bằng ngôn ngữ tự nhiên đã được đưa đến tận tay hàng trăm triệu người dùng trên toàn cầu. Khả năng "Zero-shot" và "Few-shot learning" - giải quyết các bài toán mới mà không cần phải tinh chỉnh (fine-tuning) mô hình - đã cách mạng hóa cách chúng ta tương tác với máy tính.

Không chỉ dừng lại ở văn bản, AI hiện đại có khả năng đa phương thức (multimodal). Các hệ thống như Midjourney, DALL-E 3 hay Stable Diffusion có thể tạo ra các tác phẩm nghệ thuật, nhiếp ảnh có độ phân giải siêu cao và độ chân thực cực lớn chỉ từ một câu mô tả (prompt). Trong lĩnh vực y sinh học, hệ thống AlphaFold của DeepMind đã giải quyết thành công bài toán dự đoán cấu trúc protein 3D, mở ra kỷ nguyên mới cho ngành phát hiện và điều chế thuốc (drug discovery).

Mặc dù đạt được những thành tựu phi thường, Trí tuệ nhân tạo hiện tại không phải là không có giới hạn và đối mặt với nhiều thách thức sâu sắc:

\textbf{1. Ảo giác (Hallucination):} Các mô hình ngôn ngữ lớn bản chất là những bộ dự đoán xác suất vô cùng phức tạp, không phải là cơ sở dữ liệu tri thức tĩnh. Chúng có thể tự tin sinh ra các thông tin hoàn toàn sai lệch hoặc bịa đặt các nguồn tài liệu khoa học không tồn tại.

\textbf{2. Khả năng diễn giải (Explainable AI - XAI):} Các mạng nơ-ron sâu với hàng trăm tỷ tham số hoạt động như những "hộp đen" (black box). Các nhà nghiên cứu thường không thể hiểu được cơ chế tại sao một mô hình đưa ra một quyết định hoặc một đầu ra cụ thể. Điều này là một rào cản chí mạng khi áp dụng AI trong các lĩnh vực có tính rủi ro cao như y tế, luật pháp, và xe tự lái, nơi trách nhiệm giải trình (accountability) là bắt buộc.

\textbf{3. Thiên kiến và Đạo đức (Bias and Ethics):} Vì AI được huấn luyện trên dữ liệu lấy từ thế giới thực (vốn luôn chứa đựng những định kiến về chủng tộc, giới tính, tôn giáo), các mô hình thường phản ánh và thậm chí khuếch đại các thiên kiến (bias) này, dẫn đến những phân biệt đối xử tự động trong tuyển dụng hoặc xét duyệt hồ sơ tín dụng. Các vấn đề bản quyền dữ liệu liên quan đến nội dung sáng tạo cũng đang được tranh luận pháp lý rất căng thẳng trên toàn cầu.

\textbf{4. Chi phí môi trường và Phần cứng:} Việc huấn luyện các LLM lớn tốn kém hàng chục triệu đô la và tiêu thụ một lượng điện năng khổng lồ, đi kèm với lượng khí thải carbon cao. Việc tối ưu hóa (quantization, pruning) và phát triển phần cứng xanh, hiệu quả hơn là một hướng nghiên cứu cực kỳ cấp bách.

Tương lai của AI hứa hẹn hướng tới những mô hình suy luận tốt hơn, có thể mô phỏng một phần khả năng lên kế hoạch dài hạn (Agentic AI) và khả năng suy luận logic dựa trên toán học. Cuộc chạy đua công nghệ đang diễn ra gay gắt để tìm kiếm "Trí tuệ nhân tạo tổng quát" (Artificial General Intelligence - AGI) - loại AI có khả năng hiểu, học hỏi và áp dụng kiến thức trong mọi lĩnh vực nhận thức ở mức ngang bằng hoặc vượt trội so với trí tuệ con người. Liệu AGI là một ước mơ kiến tạo nên một kỷ nguyên thịnh vượng mới cho nhân loại hay là một rủi ro hiện sinh (existential risk), như nhận định của nhiều triết gia và nhà khoa học hàng đầu, vẫn là câu hỏi lớn nhất của thế kỷ 21.

\vspace{1cm}
\textbf{Tổng kết:}
Trí tuệ nhân tạo đã đi một chặng đường dài từ những phỏng đoán lý thuyết của Alan Turing và các thuật toán logic của Hội nghị Dartmouth, vượt qua những "mùa đông" khó khăn về kỹ thuật và sự hoài nghi của xã hội, để vươn tới cuộc bùng nổ của sức mạnh tính toán hiện đại và Học sâu. Từ một lĩnh vực nghiên cứu học thuật biệt lập, AI đã trở thành một cơ sở hạ tầng nền tảng, một động lực chính đang tái cấu trúc nền kinh tế, khoa học và xã hội loài người. Chương tiếp theo sẽ đi sâu vào các cấu trúc thuật toán nền tảng tạo nên sức mạnh cho các hệ thống Deep Learning đương đại.

\section{Các kỹ thuật Machine Learning cơ bản}
\subsection{Cơ sở toán học của Mạng nơ-ron Truyền thẳng Đa tầng (MLP)}
Mạng nơ-ron truyền thẳng đa tầng (Multi-Layer Perceptron - MLP) là một trong những kiến trúc kinh điển nhất của học sâu, được thiết kế để giải quyết các bài toán hồi quy và phân loại phi tuyến tính. Một MLP bao gồm một tầng đầu vào (Input layer), một hoặc nhiều tầng ẩn (Hidden layers), và một tầng đầu ra (Output layer).

Quá trình tính toán lan truyền xuôi (Forward Propagation) từ tầng $l-1$ sang tầng $l$ được định nghĩa bằng các phương trình đại số tuyến tính:
\begin{equation}
    Z^{(l)} = W^{(l)} \cdot A^{(l-1)} + b^{(l)}
\end{equation}
\begin{equation}
    A^{(l)} = g^{(l)}(Z^{(l)})
\end{equation}
Trong đó, $W^{(l)}$ là ma trận trọng số (weight matrix) kết nối các nơ-ron ở tầng $l-1$ tới tầng $l$, $b^{(l)}$ là vector độ chệch (bias), và $g^{(l)}(\cdot)$ là hàm kích hoạt phi tuyến tính. Việc sử dụng các hàm kích hoạt như ReLU ($g(z) = \max(0, z)$) hay Sigmoid ($g(z) = \frac{1}{1 + e^{-z}}$) là bắt buộc. Nếu loại bỏ hàm kích hoạt, mạng lưới dù có sâu đến hàng trăm tầng ẩn cũng sẽ suy biến hoàn toàn về một phép biến đổi tuyến tính đơn nhất.

\subsection{Thuật toán Tối ưu hóa và Lan truyền ngược (Backpropagation)}
Để mạng nơ-ron có thể "học" được từ dữ liệu, ta cần xác định một hàm mất mát (Loss function) $J(W, b)$ nhằm định lượng mức độ sai lệch giữa dự đoán của mô hình và giá trị thực tế. Với bài toán hồi quy, hàm Sai số Toàn phương Trung bình (Mean Squared Error - MSE) thường được sử dụng. Đối với bài toán phân loại, hàm Cross-Entropy là lựa chọn tối ưu.

Mục tiêu của quá trình huấn luyện là tìm kiếm bộ tham số $\theta = \{W^{(l)}, b^{(l)}\}$ sao cho hàm $J$ đạt giá trị nhỏ nhất. Quá trình này được thực hiện thông qua thuật toán Gradient Descent (Hạ gradient):
\begin{equation}
    W^{(l)} := W^{(l)} - \alpha \frac{\partial J}{\partial W^{(l)}}
\end{equation}
\begin{equation}
    b^{(l)} := b^{(l)} - \alpha \frac{\partial J}{\partial b^{(l)}}
\end{equation}
với $\alpha$ là tốc độ học (learning rate).

Để tính toán các đạo hàm riêng $\frac{\partial J}{\partial W}$ và $\frac{\partial J}{\partial b}$ ở các tầng ẩn sâu bên trong, thuật toán Lan truyền ngược áp dụng quy tắc chuỗi (Chain rule) của giải tích. Sai số $\delta^{(l)}$ tại tầng $l$ được tính ngược từ sai số $\delta^{(l+1)}$ của tầng $l+1$:
\begin{equation}
    \delta^{(l)} = (W^{(l+1)})^T \delta^{(l+1)} \odot (g^{(l)})'(Z^{(l)})
\end{equation}
Từ đó, gradient của trọng số và bias được tính toán:
\begin{equation}
    \frac{\partial J}{\partial W^{(l)}} = \delta^{(l)} (A^{(l-1)})^T, \quad \frac{\partial J}{\partial b^{(l)}} = \delta^{(l)}
\end{equation}

\subsection{Thực nghiệm 1: Bộ dữ liệu Y tế Diabetes Pima}
Bộ dữ liệu Diabetes từ Kaggle bao gồm 768 bệnh nhân nữ, với 8 đặc trưng y tế (như nồng độ Glucose, Huyết áp, Chỉ số BMI, Tuổi tác) nhằm mục đích dự đoán khả năng mắc bệnh tiểu đường (Phân loại nhị phân).
Quá trình phân tích dữ liệu thăm dò (EDA) cho thấy bộ dữ liệu có sự chênh lệch lớn về thang đo giữa các biến. Ví dụ, mức độ Insulin dao động từ 0 đến 846, trong khi chỉ số DiabetesPedigreeFunction chỉ nằm trong khoảng 0.07 đến 2.42. Do đó, kỹ thuật `StandardScaler` (Z-score normalization) đã được áp dụng để chuẩn hóa mọi đặc trưng về phân phối có trung bình bằng 0 và phương sai bằng 1.

\subsubsection{Cài đặt thuật toán từ đầu bằng NumPy (From Scratch)}
Để hiểu rõ cơ chế nội tại, mô hình được lập trình thủ công hoàn toàn bằng NumPy. Quá trình forward pass và tính toán đạo hàm backward pass được triển khai minh bạch:
\begin{lstlisting}[language=Python]
class MLP_NumPy:
    def __init__(self, input_dim):
        self.W1 = np.random.randn(input_dim, 16) * np.sqrt(2./input_dim) # He initialization
        self.b1 = np.zeros((1, 16))
        self.W2 = np.random.randn(16, 1) * 0.01
        self.b2 = np.zeros((1, 1))

    def forward(self, X):
        self.Z1 = np.dot(X, self.W1) + self.b1
        self.A1 = np.maximum(0, self.Z1) # ReLU activation
        self.Z2 = np.dot(self.A1, self.W2) + self.b2
        self.A2 = 1 / (1 + np.exp(-self.Z2)) # Sigmoid activation
        return self.A2
        
    def backward(self, X, y, lr=0.01):
        m = X.shape[0]
        dZ2 = self.A2 - y
        dW2 = (1/m) * np.dot(self.A1.T, dZ2)
        db2 = (1/m) * np.sum(dZ2, axis=0, keepdims=True)
        dZ1 = np.dot(dZ2, self.W2.T) * (self.Z1 > 0) # ReLU derivative
        dW1 = (1/m) * np.dot(X.T, dZ1)
        db1 = (1/m) * np.sum(dZ1, axis=0, keepdims=True)
        
        self.W1 -= lr * dW1
        self.b1 -= lr * db1
        self.W2 -= lr * dW2
        self.b2 -= lr * db2
\end{lstlisting}

\subsubsection{Cài đặt bằng Keras và PyTorch}
Sử dụng Keras `Sequential` API, việc xây dựng mô hình được rút gọn chỉ còn vài dòng code với lớp `Dense`. Hàm mất mát `binary_crossentropy` và thuật toán tối ưu Adam được sử dụng. Trong khi đó, PyTorch sử dụng `nn.Module` kết hợp với `BCEWithLogitsLoss()` để tính toán loss trực tiếp từ logits, giúp tối ưu hóa sự ổn định số học so với việc chạy qua hàm Sigmoid trước. Cả ba framework đều đạt độ chính xác (Accuracy) trên tập kiểm thử ở mức $\sim 78\%$, nhưng tốc độ huấn luyện của Keras trên dữ liệu nhỏ bị giới hạn bởi thời gian khởi tạo đồ thị tính toán tĩnh.

\subsection{Thực nghiệm 2: Bộ dữ liệu Bất động sản California Housing}
Bộ dữ liệu California Housing là một bài toán Hồi quy (Regression) kinh điển, bao gồm thông tin về các khối nhà (block groups) tại California với hơn 20.000 bản ghi. Mục tiêu là dự báo Giá nhà trung bình (`median_house_value`) dựa trên các đặc trưng như vĩ độ, kinh độ, tổng số phòng, số người ở và thu nhập trung bình.

Dữ liệu chứa nhiều giá trị ngoại lai (outliers). Hàm mục tiêu được cấu hình là Sai số Toàn phương Trung bình (Mean Squared Error - MSE). Thay vì hàm Sigmoid ở lớp đầu ra, mô hình kết thúc bằng một lớp Linear Activation nhằm cho phép dự đoán các số thực tự do không bị giới hạn trong khoảng $[0, 1]$. Trên PyTorch, hàm `nn.MSELoss()` được áp dụng cùng với thuật toán Adam (tốc độ học $0.01$). Kết quả đối sánh cho thấy mô hình MLP 2 tầng ẩn vượt qua khả năng dự đoán của Linear Regression truyền thống tới 25\% về mặt chỉ số R-squared.

\subsection{Thực nghiệm 3: Đánh giá Phản hồi Khách hàng E-Commerce}
Bộ dữ liệu "Women's E-Commerce Clothing Reviews" là dữ liệu thực tế từ các đánh giá mua sắm. Ở bài toán này, đặc trưng dạng văn bản tự do (Review Text) được trích xuất bằng thuật toán TF-IDF (`TfidfVectorizer` trong scikit-learn). Dữ liệu dạng thưa (Sparse Matrix) với kích thước hàng chục ngàn chiều được truyền vào mô hình nơ-ron để phân loại liệu khách hàng có giới thiệu sản phẩm hay không (Recommended IND).

Kiến trúc mạng được tinh chỉnh đặc biệt với kỹ thuật Dropout ($p=0.5$) để chống lại hiện tượng quá khớp (Overfitting) vốn rất dễ xảy ra khi số lượng đặc trưng văn bản TF-IDF lớn hơn số lượng mẫu huấn luyện. Kết quả đo lường bằng ma trận nhầm lẫn (Confusion Matrix) và F1-Score minh chứng sức mạnh của Deep Learning so với các kỹ thuật thống kê cổ điển.

\section{Mạng Nơ-ron Tích chập (CNN)}
\subsection{Lý thuyết cốt lõi của Tích chập (Convolution)}
Mạng nơ-ron tích chập (CNN) là một lớp đặc biệt của mạng học sâu, được tinh chỉnh hoàn hảo để xử lý cấu trúc dữ liệu dạng lưới, đặc biệt là hình ảnh 2 chiều. Khác với MLP - nơi mọi nơ-ron ở một tầng đều kết nối chằng chịt tới toàn bộ nơ-ron ở tầng sau (Fully Connected) dẫn tới sự bùng nổ tham số - CNN áp dụng triết lý Trọng số Chia sẻ (Weight Sharing) và Kết nối Cục bộ (Local Connectivity).

Linh hồn của CNN là phép Tích chập (Convolution). Thay vì sử dụng ma trận trọng số, CNN sử dụng các Bộ lọc trượt (Kernels/Filters). Một kernel là một ma trận kích thước nhỏ (ví dụ $3 \times 3$ hoặc $5 \times 5$). Trong quá trình Forward Pass, kernel này sẽ trượt tuần tự qua toàn bộ bức ảnh từ trái sang phải, từ trên xuống dưới. Tại mỗi vị trí, một phép tính tổng các tích vô hướng (dot product) giữa kernel và phần ảnh nằm bên dưới nó được thực hiện, tạo ra một Bản đồ Đặc trưng (Feature Map). 

Về mặt toán học liên tục, tích chập 2 chiều giữa một hàm hình ảnh $I$ và một kernel $K$ được định nghĩa là:
\begin{equation}
    (I * K)(i, j) = \sum_{m} \sum_{n} I(i-m, j-n) K(m, n)
\end{equation}
Trong thực tế lập trình máy tính, phép toán này thường được triển khai dưới dạng Cross-Correlation:
\begin{equation}
    S(i, j) = \sum_{m=0}^{k-1} \sum_{n=0}^{k-1} I(i \cdot s + m, j \cdot s + n) K(m, n)
\end{equation}
với $s$ là bước sải (stride).

Ngoài lớp Convolution, mạng CNN luôn đi kèm với lớp Gộp (Pooling Layer), phổ biến nhất là Max Pooling. Nhiệm vụ của lớp này là hạ lấy mẫu (downsampling) dữ liệu, giảm thiểu kích thước không gian của Feature Map, từ đó giảm số lượng tham số tính toán và tạo ra tính chất Bất biến Dịch chuyển (Translation Invariance) - giúp mô hình nhận diện được vật thể dù nó nằm ở bất kỳ đâu trong khung hình.

\subsection{Khảo sát Dataset 1: Nhận diện chữ số viết tay MNIST}
MNIST (Modified National Institute of Standards and Technology) là bộ dữ liệu "Hello World" của Computer Vision. Nó chứa 60.000 ảnh huấn luyện và 10.000 ảnh kiểm thử, mỗi ảnh là một chữ số viết tay từ 0 đến 9, kích thước $28 \times 28$ pixel ở định dạng xám (grayscale).

Kiến trúc mạng được thiết kế tuân theo phong cách LeNet-5 cổ điển. Khởi đầu bằng lớp `Conv2D` với 16 bộ lọc kích thước $3 \times 3$, tiếp nối bằng một lớp `MaxPooling2D` ($2 \times 2$), sau đó là một lớp `Flatten` để trải phẳng Bản đồ đặc trưng 2D thành vector 1D, và cuối cùng là lớp Dense kết nối đầy đủ (Fully Connected) phân loại 10 lớp thông qua hàm kích hoạt Softmax.

Khi triển khai trực tiếp bằng mã giả NumPy, phép trượt Convolution đòi hỏi 4 vòng lặp for lồng nhau (vòng lặp Batch, vòng lặp Kênh, vòng lặp Chiều cao và Chiều rộng), gây ra độ trễ điện toán khổng lồ. Tuy nhiên, khi chuyển sang Keras và PyTorch, kiến trúc đồ thị tính toán biên dịch trực tiếp phép Convolution thành các lời gọi hàm cuDNN trên card đồ họa NVIDIA. Kết quả huấn luyện sau 10 Epochs với Keras đạt độ chính xác kiểm thử trên $99.2\%$.

\subsection{Khảo sát Dataset 2: Nhận diện vật thể màu CIFAR-10}
CIFAR-10 phức tạp hơn nhiều so với MNIST. Bộ dữ liệu chứa 60.000 ảnh màu độ phân giải $32 \times 32$, đại diện cho 10 lớp vật thể (Máy bay, Ô tô, Chim, Mèo, Hươu, Chó, Ếch, Ngựa, Tàu thủy, Xe tải). Sự phức tạp nằm ở chỗ hình ảnh chứa nhiễu bối cảnh, góc chụp đa dạng, và có 3 kênh màu (Red, Green, Blue) tạo thành một tensor đầu vào có kích thước $[Batch\_Size, 32, 32, 3]$.

Để đối phó với mức độ phức tạp này, một kiến trúc VGG-style thu nhỏ đã được xây dựng. Các lớp `Conv2D` được xếp chồng lên nhau liên tiếp trước khi gọi lớp `MaxPooling2D`. Việc này cho phép mạng học được các đặc trưng phi tuyến sâu hơn (nhờ hàm ReLU chèn vào giữa) mà không làm suy giảm ngay lập tức thông tin không gian. Thư viện `torchvision.datasets` của PyTorch được sử dụng để tự động tải và định dạng Tensor. Hàm mất mát `CrossEntropyLoss` (đã bao gồm log-softmax) được dùng để huấn luyện. 

\subsection{Khảo sát Dataset 3: Thực nghiệm Tích chập 1 chiều (1D-CNN) trên Dữ liệu Bảng Diabetes}
Thông thường, dữ liệu dạng bảng (Tabular Data) được giải quyết bằng cấu trúc MLP, Random Forest, hoặc XGBoost. Tuy nhiên, trong nghiên cứu này, một thử nghiệm đột phá đã được tiến hành: ứng dụng CNN 1 chiều (Conv1D) để trích xuất đặc trưng của tập dữ liệu y tế Diabetes.

Dữ liệu chứa 21 đặc trưng (features) của từng bệnh nhân được coi như một "hình ảnh" có chiều cao 21 pixel và chiều rộng 1 pixel. Lớp `Conv1D` với cửa sổ trượt chiều dọc $K=3$ được lướt qua các biến y tế này. Ý tưởng cốt lõi là mạng nơ-ron có thể khám phá ra các mối tương quan cục bộ liền kề giữa các nhóm biến lân cận. Sau lớp Conv1D, dữ liệu đi qua lớp `GlobalMaxPooling1D` để trích xuất tín hiệu mạnh nhất trên toàn bộ trục không gian đặc trưng, rồi đưa vào bộ phân loại. Mô hình Keras cho bài toán Conv1D này chỉ mất khoảng 2 giây/epoch để hội tụ, minh chứng cho tính gọn nhẹ và linh hoạt của mạng tích chập vượt ra ngoài giới hạn xử lý ảnh không gian 2D thông thường.

\section{Mạng Nơ-ron Hồi quy (RNN) và Xử lý Dữ liệu Chuỗi Thời gian}

Trong chương này, chúng ta sẽ đi sâu vào việc nghiên cứu Mạng Nơ-ron Hồi quy (Recurrent Neural Networks - RNN) và các biến thể mạnh mẽ của nó bao gồm Bộ nhớ Ngắn-Dài hạn (Long Short-Term Memory - LSTM) và Cổng Hồi quy (Gated Recurrent Unit - GRU). Kiến trúc RNN được thiết kế đặc biệt để xử lý dữ liệu dạng chuỗi, trong đó thông tin từ các bước thời gian trước đó được lưu giữ và sử dụng để dự đoán hoặc phân loại tại các bước thời gian hiện tại.

\subsection{Cơ sở Lý thuyết cơ bản}

\subsubsection{Mạng Nơ-ron Hồi quy (RNN)}
Mạng Nơ-ron Hồi quy cơ bản duy trì một trạng thái ẩn (hidden state) $h_t$ tại mỗi bước thời gian $t$. Trạng thái này đóng vai trò như một "bộ nhớ" lưu trữ thông tin về các phần tử đã quan sát trước đó trong chuỗi.
Công thức tính toán của một mạng RNN cơ bản được định nghĩa như sau:
\begin{align}
    h_t &= \tanh(W_{hh} h_{t-1} + W_{xh} x_t + b_h) \\
    y_t &= W_{hy} h_t + b_y
\end{align}
Trong đó:
\begin{itemize}
    \item $x_t \in \mathbb{R}^d$: Vector đầu vào tại thời điểm $t$.
    \item $h_t \in \mathbb{R}^h$: Trạng thái ẩn tại thời điểm $t$.
    \item $W_{hh} \in \mathbb{R}^{h \times h}$, $W_{xh} \in \mathbb{R}^{h \times d}$, $W_{hy} \in \mathbb{R}^{o \times h}$: Các ma trận trọng số.
    \item $b_h \in \mathbb{R}^h$, $b_y \in \mathbb{R}^o$: Các vector bias.
    \item $\tanh$: Hàm kích hoạt phi tuyến (Hyperbolic Tangent).
\end{itemize}

\subsubsection{Lan truyền ngược qua thời gian (BPTT)}
Để huấn luyện mạng RNN, chúng ta sử dụng thuật toán Lan truyền ngược qua thời gian (Backpropagation Through Time - BPTT). Quá trình này về cơ bản là việc mở ra (unroll) mạng RNN theo các bước thời gian và áp dụng thuật toán lan truyền ngược tiêu chuẩn. Tuy nhiên, đạo hàm của hàm mất mát $L$ theo trọng số $W_{hh}$ liên quan đến việc nhân liên tiếp các ma trận đạo hàm riêng của trạng thái ẩn:
\begin{equation}
    \frac{\partial L}{\partial W_{hh}} = \sum_{t=1}^T \sum_{k=1}^t \frac{\partial L_t}{\partial y_t} \frac{\partial y_t}{\partial h_t} \left( \prod_{j=k+1}^t \frac{\partial h_j}{\partial h_{j-1}} \right) \frac{\partial h_k}{\partial W_{hh}}
\end{equation}
Nếu các giá trị riêng của ma trận $\frac{\partial h_j}{\partial h_{j-1}}$ nhỏ hơn 1, quá trình nhân liên tiếp sẽ làm gradient tiến dần về 0 (Vanishing Gradient). Ngược lại, nếu lớn hơn 1, gradient sẽ bùng nổ (Exploding Gradient). Điều này khiến mạng RNN cơ bản khó học được các phụ thuộc xa trong chuỗi (long-term dependencies).

\subsubsection{Bộ nhớ Ngắn-Dài hạn (LSTM)}
Để giải quyết bài toán Vanishing Gradient, kiến trúc LSTM được đề xuất với một cấu trúc ô nhớ (cell state) $c_t$ và ba cổng điều khiển luồng thông tin: cổng quên (forget gate), cổng đầu vào (input gate), và cổng đầu ra (output gate).
\begin{align}
    f_t &= \sigma(W_f \cdot [h_{t-1}, x_t] + b_f) \quad \text{(Cổng quên)} \\
    i_t &= \sigma(W_i \cdot [h_{t-1}, x_t] + b_i) \quad \text{(Cổng đầu vào)} \\
    \tilde{c}_t &= \tanh(W_c \cdot [h_{t-1}, x_t] + b_c) \quad \text{(Trạng thái ô nhớ ứng viên)} \\
    c_t &= f_t \odot c_{t-1} + i_t \odot \tilde{c}_t \quad \text{(Cập nhật ô nhớ)} \\
    o_t &= \sigma(W_o \cdot [h_{t-1}, x_t] + b_o) \quad \text{(Cổng đầu ra)} \\
    h_t &= o_t \odot \tanh(c_t) \quad \text{(Trạng thái ẩn mới)}
\end{align}
Nhờ phép cộng tuyến tính tại bước cập nhật ô nhớ $c_t$, gradient có thể luân chuyển dễ dàng qua các bước thời gian mà không bị suy giảm nhanh chóng.

\subsubsection{Gated Recurrent Unit (GRU)}
GRU là một biến thể rút gọn của LSTM, kết hợp cổng quên và cổng đầu vào thành một "cổng cập nhật" (update gate) $z_t$, đồng thời sử dụng "cổng thiết lập lại" (reset gate) $r_t$ thay vì có cell state tách biệt:
\begin{align}
    z_t &= \sigma(W_z \cdot [h_{t-1}, x_t] + b_z) \\
    r_t &= \sigma(W_r \cdot [h_{t-1}, x_t] + b_r) \\
    \tilde{h}_t &= \tanh(W \cdot [r_t \odot h_{t-1}, x_t] + b) \\
    h_t &= (1 - z_t) \odot h_{t-1} + z_t \odot \tilde{h}_t
\end{align}
GRU hoạt động nhanh hơn và tốn ít bộ nhớ hơn LSTM trong khi thường đạt được độ chính xác tương đương ở nhiều bài toán dự báo chuỗi thời gian.

\newpage
% =========================================================================================
% DATASET 1: AMAZON
% =========================================================================================
\subsection{Tập dữ liệu 1: Dự báo Giá Cổ phiếu Amazon (AMZN)}

\subsubsection{Mô tả, Phân phối và Phân tích Dữ liệu}
Tập dữ liệu giá cổ phiếu Amazon (AMZN Stock Price) chứa thông tin giao dịch hàng ngày trên thị trường chứng khoán. 
\begin{itemize}
    \item \textbf{Link}: Tải về từ Yahoo Finance / Kaggle.
    \item \textbf{Kích thước}: Hơn 5000 dòng, bao gồm các cột như Date, Open, High, Low, Close, Adj Close, Volume.
    \item \textbf{Sample Outputs}: Tại ngày giao dịch, ví dụ "2023-01-10", giá trị Close có thể là 89.87 USD.
    \item \textbf{Phân phối}: Dữ liệu có xu hướng tăng mạnh từ năm 2010 đến 2021, đạt đỉnh, sau đó có sự điều chỉnh. Đồ thị giá đóng cửa có tính chất không dừng (non-stationary) rõ rệt, đặc trưng của chuỗi tài chính. Phân phối của sai phân (returns) tiệm cận phân phối chuẩn nhưng có đuôi dày (fat tails).
    \item \textbf{Phân tích (EDA)}: Tính tự tương quan (Autocorrelation) cho thấy giá ngày hôm nay phụ thuộc cực kỳ chặt chẽ vào ngày trước đó. Việc dự đoán yêu cầu trích xuất xu hướng ngắn hạn và loại bỏ nhiễu.
\end{itemize}

\subsubsection{Cài đặt mô hình bằng Python (From Scratch - NumPy)}
Để minh họa rõ nguyên lý hoạt động, đoạn mã sau trình bày cách cài đặt một mô hình RNN cơ bản bằng NumPy hoàn toàn từ đầu (From Scratch).

\begin{lstlisting}[language=Python, caption=Cài đặt RNN từ đầu với NumPy cho tập dữ liệu AMZN]
import numpy as np
import pandas as pd
import matplotlib.pyplot as plt

# 1. Chuan bi du lieu
def prepare_data(series, seq_len):
    X, y = [], []
    for i in range(len(series) - seq_len):
        X.append(series[i:i+seq_len])
        y.append(series[i+seq_len])
    return np.array(X), np.array(y)

# Gia su du lieu da duoc load va chuan hoa MinMax
# data = pd.read_csv('AMZN.csv')['Close'].values
# norm_data = (data - np.min(data)) / (np.max(data) - np.min(data))
# X_train, y_train = prepare_data(norm_data, seq_len=60)

class SimpleRNN:
    def __init__(self, input_size, hidden_size, output_size):
        # Khoi tao trong so (Xavier/He initialization)
        self.W_xh = np.random.randn(hidden_size, input_size) * 0.01
        self.W_hh = np.random.randn(hidden_size, hidden_size) * 0.01
        self.W_hy = np.random.randn(output_size, hidden_size) * 0.01
        self.b_h = np.zeros((hidden_size, 1))
        self.b_y = np.zeros((output_size, 1))
        
    def forward(self, inputs):
        # inputs shape: (seq_len, input_size)
        h_states, x_states = {}, {}
        h_states[-1] = np.zeros((self.W_hh.shape[0], 1))
        
        for t in range(len(inputs)):
            x_t = inputs[t].reshape(-1, 1)
            # h_t = tanh(W_xh * x_t + W_hh * h_{t-1} + b_h)
            h_states[t] = np.tanh(np.dot(self.W_xh, x_t) + np.dot(self.W_hh, h_states[t-1]) + self.b_h)
            x_states[t] = x_t
            
        # y = W_hy * h_last + b_y
        y_pred = np.dot(self.W_hy, h_states[len(inputs)-1]) + self.b_y
        return y_pred, h_states, x_states
        
    def backward(self, y_pred, y_true, h_states, x_states, learning_rate=0.01):
        # BPTT
        seq_len = len(x_states)
        dy = y_pred - y_true
        
        dW_hy = np.dot(dy, h_states[seq_len-1].T)
        db_y = dy
        
        dh = np.dot(self.W_hy.T, dy)
        dW_hh = np.zeros_like(self.W_hh)
        dW_xh = np.zeros_like(self.W_xh)
        db_h = np.zeros_like(self.b_h)
        
        # Lan truyen nguoc qua tung buoc thoi gian
        for t in reversed(range(seq_len)):
            temp = (1 - h_states[t] ** 2) * dh # dao ham cua tanh
            db_h += temp
            dW_xh += np.dot(temp, x_states[t].T)
            dW_hh += np.dot(temp, h_states[t-1].T)
            dh = np.dot(self.W_hh.T, temp)
            
        # Cap nhat (Gradient Descent)
        # Ap dung Gradient Clipping de tranh Exploding Gradients
        for d in [dW_xh, dW_hh, dW_hy, db_h, db_y]:
            np.clip(d, -5, 5, out=d)
            
        self.W_xh -= learning_rate * dW_xh
        self.W_hh -= learning_rate * dW_hh
        self.W_hy -= learning_rate * dW_hy
        self.b_h -= learning_rate * db_h
        self.b_y -= learning_rate * db_y
        
# Vong lap huan luyen mau (Pseudo-code)
# rnn = SimpleRNN(input_size=1, hidden_size=64, output_size=1)
# for epoch in range(100):
#     loss = 0
#     for i in range(len(X_train)):
#         y_pred, h_states, x_states = rnn.forward(X_train[i])
#         loss += (y_pred - y_train[i])**2
#         rnn.backward(y_pred, y_train[i], h_states, x_states)
#     print(f"Epoch {epoch}, Loss: {loss/len(X_train)}")
\end{lstlisting}

\subsubsection{Cài đặt mô hình bằng Keras}
Keras cung cấp các API bậc cao giúp việc xây dựng các mô hình chuỗi thời gian phức tạp trở nên dễ dàng hơn. Dưới đây là kiến trúc LSTM kết hợp với Dropout để chống Overfitting cho dữ liệu AMZN.

\begin{lstlisting}[language=Python, caption=Cài đặt LSTM bằng Keras cho dữ liệu AMZN]
import tensorflow as tf
from tensorflow.keras.models import Sequential
from tensorflow.keras.layers import LSTM, Dense, Dropout

def build_keras_lstm(seq_len, input_dim):
    model = Sequential()
    
    # Lop LSTM thu nhat
    model.add(LSTM(units=100, return_sequences=True, input_shape=(seq_len, input_dim)))
    model.add(Dropout(0.2))
    
    # Lop LSTM thu hai
    model.add(LSTM(units=100, return_sequences=False))
    model.add(Dropout(0.2))
    
    # Lop dau ra (Du doan gia tri lien tuc - Regression)
    model.add(Dense(units=25))
    model.add(Dense(units=1))
    
    # Compile model voi Adam optimizer va huber loss
    model.compile(optimizer='adam', loss='mean_squared_error', metrics=['mae'])
    
    return model

# Thuc thi huan luyen
# model_keras = build_keras_lstm(seq_len=60, input_dim=1)
# history = model_keras.fit(X_train, y_train, batch_size=32, epochs=50, validation_split=0.1)
\end{lstlisting}

\subsubsection{Cài đặt mô hình bằng PyTorch}
Đối với những nhà nghiên cứu cần sự tuỳ biến cao về vòng lặp huấn luyện, PyTorch là một lựa chọn tuyệt vời. Đoạn mã dưới đây sử dụng GRU cho tập AMZN.

\begin{lstlisting}[language=Python, caption=Cài đặt GRU bằng PyTorch cho dữ liệu AMZN]
import torch
import torch.nn as nn
from torch.utils.data import TensorDataset, DataLoader

class PyTorchGRU(nn.Module):
    def __init__(self, input_dim, hidden_dim, num_layers, output_dim):
        super(PyTorchGRU, self).__init__()
        self.hidden_dim = hidden_dim
        self.num_layers = num_layers
        
        # Dinh nghia lop GRU
        self.gru = nn.GRU(input_dim, hidden_dim, num_layers, batch_first=True, dropout=0.2)
        
        # Dinh nghia lop Fully Connected
        self.fc = nn.Linear(hidden_dim, output_dim)
        
    def forward(self, x):
        # Khoi tao hidden state voi cac gia tri 0
        h0 = torch.zeros(self.num_layers, x.size(0), self.hidden_dim).requires_grad_()
        
        # Chuyen tiep qua GRU
        out, (hn) = self.gru(x, (h0.detach()))
        
        # Lay dau ra o buoc thoi gian cuoi cung
        out = self.fc(out[:, -1, :]) 
        return out

# Huan luyen mo hinh
# model_pt = PyTorchGRU(input_dim=1, hidden_dim=64, num_layers=2, output_dim=1)
# criterion = nn.MSELoss(reduction='mean')
# optimizer = torch.optim.Adam(model_pt.parameters(), lr=0.001)

# X_train_t = torch.Tensor(X_train).unsqueeze(-1)
# y_train_t = torch.Tensor(y_train).unsqueeze(-1)
# loader = DataLoader(TensorDataset(X_train_t, y_train_t), batch_size=64, shuffle=True)

# for epoch in range(50):
#     for inputs, targets in loader:
#         optimizer.zero_grad()
#         outputs = model_pt(inputs)
#         loss = criterion(outputs, targets)
#         loss.backward()
#         optimizer.step()
\end{lstlisting}

\newpage
% =========================================================================================
% DATASET 2: GOLD PRICE
% =========================================================================================
\subsection{Tập dữ liệu 2: Dự báo Giá Vàng (Gold Price)}

\subsubsection{Mô tả, Phân phối và Phân tích Dữ liệu}
Vàng là một trong những tài sản trú ẩn an toàn phổ biến nhất, chịu ảnh hưởng từ nhiều yếu tố vĩ mô như lạm phát, lãi suất, và biến động địa chính trị.
\begin{itemize}
    \item \textbf{Link}: Dữ liệu giao dịch vàng giao ngay (XAU/USD) từ Investing.com.
    \item \textbf{Kích thước}: Khoảng 4000 mẫu dữ liệu lịch sử hàng ngày.
    \item \textbf{Sample Outputs}: Open: 1850.50, High: 1870.00, Low: 1845.20, Close: 1865.30.
    \item \textbf{Phân phối}: Không giống như cổ phiếu công nghệ, giá vàng thường biến động mạnh trong các đợt khủng hoảng và đi ngang trong thời kỳ ổn định. Phân phối của lợi suất (returns) vàng hẹp hơn nhưng có độ nhọn kurtosis cao.
    \item \textbf{Phân tích (EDA)}: Tính thời vụ (seasonality) không quá rõ rệt nhưng tồn tại các chu kỳ dài hạn (long-term cycles). Điều này gợi ý rằng các kiến trúc có khả năng bắt được chuỗi thời gian dài như LSTM hay GRU sẽ rất phù hợp.
\end{itemize}

\subsubsection{Cài đặt mô hình bằng Python (From Scratch - NumPy)}
Trong phần này, ta cài đặt một phiên bản chi tiết của LSTM tiến triển từ đầu.

\begin{lstlisting}[language=Python, caption=Cài đặt thành phần LSTM từ đầu với NumPy cho Giá Vàng]
import numpy as np

def sigmoid(x):
    return 1 / (1 + np.exp(-np.clip(x, -500, 500)))

class ScratchLSTM:
    def __init__(self, input_dim, hidden_dim):
        self.hidden_dim = hidden_dim
        # Trong so cho cac cong: Input, Forget, Output, Cell
        # Gop chung vao mot ma tran lon de toi uu hoa
        self.W_f = np.random.randn(hidden_dim, hidden_dim + input_dim) * 0.1
        self.b_f = np.zeros((hidden_dim, 1))
        
        self.W_i = np.random.randn(hidden_dim, hidden_dim + input_dim) * 0.1
        self.b_i = np.zeros((hidden_dim, 1))
        
        self.W_c = np.random.randn(hidden_dim, hidden_dim + input_dim) * 0.1
        self.b_c = np.zeros((hidden_dim, 1))
        
        self.W_o = np.random.randn(hidden_dim, hidden_dim + input_dim) * 0.1
        self.b_o = np.zeros((hidden_dim, 1))
        
        # Trong so dau ra
        self.W_y = np.random.randn(1, hidden_dim) * 0.1
        self.b_y = np.zeros((1, 1))
        
    def step_forward(self, x, h_prev, c_prev):
        # Noi h_prev va x
        concat_input = np.vstack((h_prev, x))
        
        # Tinh cac cong (gates)
        f_t = sigmoid(np.dot(self.W_f, concat_input) + self.b_f)
        i_t = sigmoid(np.dot(self.W_i, concat_input) + self.b_i)
        c_tilde = np.tanh(np.dot(self.W_c, concat_input) + self.b_c)
        
        # Cap nhat cell state
        c_next = f_t * c_prev + i_t * c_tilde
        
        # Tinh output gate va hidden state
        o_t = sigmoid(np.dot(self.W_o, concat_input) + self.b_o)
        h_next = o_t * np.tanh(c_next)
        
        return h_next, c_next, f_t, i_t, c_tilde, o_t, concat_input

    def forward(self, x_seq):
        seq_len = x_seq.shape[0]
        h, c = np.zeros((self.hidden_dim, 1)), np.zeros((self.hidden_dim, 1))
        caches = []
        
        for t in range(seq_len):
            x_t = x_seq[t].reshape(-1, 1)
            h, c, f, i, c_t, o, concat = self.step_forward(x_t, h, c)
            caches.append((h, c, f, i, c_t, o, concat))
            
        y_pred = np.dot(self.W_y, h) + self.b_y
        return y_pred, caches
        
    # (Phan Backward Propagation cho LSTM rat dai va phuc tap, 
    # doi hoi ap dung Chain Rule cho cac phep nhan Hadamard cua 4 cong)
\end{lstlisting}

\subsubsection{Cài đặt mô hình bằng Keras}
Mô hình dành cho giá Vàng có thể được cấu hình với RNN tiêu chuẩn hoặc GRU. Dưới đây là việc sử dụng cấu trúc RNN nhiều lớp với Keras.

\begin{lstlisting}[language=Python, caption=Cài đặt RNN cơ bản qua Keras cho Giá Vàng]
from tensorflow.keras.models import Sequential
from tensorflow.keras.layers import SimpleRNN, Dense, BatchNormalization

def build_gold_rnn(seq_len, input_dim):
    model = Sequential()
    
    # SimpleRNN nhay cam voi Gradient Vanishing, nen ta them Batch Norm
    model.add(SimpleRNN(128, return_sequences=True, input_shape=(seq_len, input_dim)))
    model.add(BatchNormalization())
    
    model.add(SimpleRNN(64, return_sequences=False))
    model.add(BatchNormalization())
    
    model.add(Dense(32, activation='relu'))
    model.add(Dense(1)) # Du doan gia tiep theo
    
    # Toi uu hoa voi thuat toan RMSprop (thuong tot cho RNN)
    model.compile(optimizer='rmsprop', loss='mse', metrics=['mae'])
    return model
\end{lstlisting}

\subsubsection{Cài đặt mô hình bằng PyTorch}
Sử dụng LSTM trong PyTorch cho mô hình Vàng, tập trung vào kiến trúc Bidirectional để xem xét tín hiệu từ hai chiều (mặc dù với bài toán chuỗi thời gian nguyên thủy, ta thường chỉ xem xét chuỗi quá khứ).

\begin{lstlisting}[language=Python, caption=Cài đặt LSTM bằng PyTorch cho Giá Vàng]
import torch.nn as nn

class GoldPyTorchLSTM(nn.Module):
    def __init__(self, input_dim, hidden_dim, num_layers, output_dim, dropout=0.3):
        super(GoldPyTorchLSTM, self).__init__()
        self.hidden_dim = hidden_dim
        self.num_layers = num_layers
        
        # Khoi tao LSTM layer
        self.lstm = nn.LSTM(input_dim, hidden_dim, num_layers, 
                            batch_first=True, dropout=dropout)
        
        # Dense layer
        self.fc1 = nn.Linear(hidden_dim, 32)
        self.relu = nn.ReLU()
        self.fc2 = nn.Linear(32, output_dim)
        
    def forward(self, x):
        h0 = torch.zeros(self.num_layers, x.size(0), self.hidden_dim).to(x.device)
        c0 = torch.zeros(self.num_layers, x.size(0), self.hidden_dim).to(x.device)
        
        # Truy truyen qua LSTM
        out, _ = self.lstm(x, (h0, c0))
        
        # Lay trang thai an tai buoc cuoi cung
        out = out[:, -1, :]
        out = self.relu(self.fc1(out))
        out = self.fc2(out)
        return out
\end{lstlisting}


\newpage
% =========================================================================================
% DATASET 3: SILVER PRICE
% =========================================================================================
\subsection{Tập dữ liệu 3: Dự báo Giá Bạc (Silver Price)}

\subsubsection{Mô tả, Phân phối và Phân tích Dữ liệu}
Giá Bạc mang tính chất song song - vừa là kim loại quý đầu tư, vừa là kim loại công nghiệp được sử dụng rộng rãi trong điện tử và năng lượng mặt trời.
\begin{itemize}
    \item \textbf{Link}: Lấy từ LBMA (London Bullion Market Association) hoặc Quandl.
    \item \textbf{Kích thước}: Khoảng 3500 dòng dữ liệu theo ngày.
    \item \textbf{Sample Outputs}: Open: 23.10, High: 23.45, Low: 22.90, Close: 23.35.
    \item \textbf{Phân phối}: Giá bạc biến động mạnh (volatility) hơn giá vàng do thị trường nhỏ hơn và bị phụ thuộc vào sản xuất công nghiệp. Điều này làm cho việc dự báo khó khăn hơn. Phân phối có dạng bimodal (hai đỉnh) tùy thuộc vào các chu kỳ bùng nổ của ngành năng lượng.
    \item \textbf{Phân tích (EDA)}: Sự tương quan chéo (Cross-correlation) với Vàng rất cao. Có hiện tượng tăng đột biến và sụp đổ nhanh, đòi hỏi mô hình phản ứng nhanh với các thay đổi ngắn hạn. GRU là cấu trúc thích hợp nhất do có khả năng học các thay đổi đột ngột hiệu quả với số lượng tham số ít hơn.
\end{itemize}

\subsubsection{Cài đặt mô hình bằng Python (From Scratch - NumPy)}
Ta minh họa kiến trúc GRU (Gated Recurrent Unit) tự viết, rất thích hợp cho dữ liệu Bạc.

\begin{lstlisting}[language=Python, caption=Cài đặt GRU từ đầu với NumPy cho Giá Bạc]
import numpy as np

def sigmoid(x): return 1 / (1 + np.exp(-x))

class ScratchGRU:
    def __init__(self, input_dim, hidden_dim):
        self.h_dim = hidden_dim
        
        # Reset gate
        self.W_r = np.random.randn(hidden_dim, hidden_dim + input_dim) * 0.1
        self.b_r = np.zeros((hidden_dim, 1))
        
        # Update gate
        self.W_z = np.random.randn(hidden_dim, hidden_dim + input_dim) * 0.1
        self.b_z = np.zeros((hidden_dim, 1))
        
        # Candidate hidden state
        self.W_h = np.random.randn(hidden_dim, hidden_dim + input_dim) * 0.1
        self.b_h = np.zeros((hidden_dim, 1))
        
        self.W_y = np.random.randn(1, hidden_dim) * 0.1
        
    def step(self, x_t, h_prev):
        concat = np.vstack((h_prev, x_t))
        
        # Calculate gates
        r_t = sigmoid(np.dot(self.W_r, concat) + self.b_r)
        z_t = sigmoid(np.dot(self.W_z, concat) + self.b_z)
        
        # Reset h_prev
        concat_reset = np.vstack((r_t * h_prev, x_t))
        h_tilde = np.tanh(np.dot(self.W_h, concat_reset) + self.b_h)
        
        # Final hidden state
        h_next = (1 - z_t) * h_prev + z_t * h_tilde
        return h_next
        
    def forward(self, x_seq):
        h = np.zeros((self.h_dim, 1))
        for t in range(x_seq.shape[0]):
            x_t = x_seq[t].reshape(-1, 1)
            h = self.step(x_t, h)
        return np.dot(self.W_y, h)
\end{lstlisting}

\subsubsection{Cài đặt mô hình bằng Keras}
Mô hình GRU đa lớp kết hợp với L1/L2 Regularization để ngăn chặn mô hình học thuộc nhiễu (noise) của giá Bạc.

\begin{lstlisting}[language=Python, caption=Cài đặt Keras GRU cho Giá Bạc]
from tensorflow.keras.models import Sequential
from tensorflow.keras.layers import GRU, Dense, Dropout
from tensorflow.keras.regularizers import l2

def build_silver_gru(seq_len, input_dim):
    model = Sequential()
    
    # Lop GRU thu nhat voi L2 regularization
    model.add(GRU(64, return_sequences=True, 
                  input_shape=(seq_len, input_dim),
                  kernel_regularizer=l2(0.001)))
    model.add(Dropout(0.3))
    
    # Lop GRU thu hai
    model.add(GRU(32, return_sequences=False))
    model.add(Dropout(0.3))
    
    # Lop ket xuat
    model.add(Dense(16, activation='relu'))
    model.add(Dense(1, activation='linear'))
    
    model.compile(optimizer='adam', loss='huber_loss', metrics=['mse'])
    return model
\end{lstlisting}

\subsubsection{Cài đặt mô hình bằng PyTorch}
Sử dụng PyTorch RNN kết hợp với Learning Rate Scheduler để tối ưu hóa hội tụ.

\begin{lstlisting}[language=Python, caption=Cài đặt PyTorch cơ bản RNN cho Giá Bạc]
import torch.nn as nn
import torch.optim as optim

class SilverPyTorchRNN(nn.Module):
    def __init__(self, input_dim, hidden_dim, num_layers, output_dim):
        super(SilverPyTorchRNN, self).__init__()
        self.rnn = nn.RNN(input_dim, hidden_dim, num_layers, 
                          batch_first=True, nonlinearity='relu')
        self.fc = nn.Linear(hidden_dim, output_dim)
        
    def forward(self, x):
        # RNN tra ve out (cac hidden state) va hn (hidden state cuoi)
        out, hn = self.rnn(x)
        # Dung hn cua lop cuoi cung de du doan
        res = self.fc(hn[-1, :, :])
        return res

# Setup mo hinh, toi uu va lap lich (Scheduler)
# model = SilverPyTorchRNN(input_dim=1, hidden_dim=64, num_layers=2, output_dim=1)
# optimizer = optim.Adam(model.parameters(), lr=0.005)
# scheduler = optim.lr_scheduler.StepLR(optimizer, step_size=20, gamma=0.5)
\end{lstlisting}

\newpage
% =========================================================================================
% EVALUATION AND COMPARISON
% =========================================================================================
\subsection{Đánh giá và So sánh Mô hình}

Để cung cấp một cái nhìn tổng quan, chúng ta áp dụng các số liệu đánh giá tiêu chuẩn như \textbf{MAE (Mean Absolute Error)} và \textbf{RMSE (Root Mean Square Error)} lên kết quả dự báo của ba kiến trúc mạng RNN cơ bản, LSTM và GRU trên từng tập dữ liệu (Đã qua chuẩn hóa MinMax). 

Dữ liệu được chia theo tỷ lệ 80\% (Training) và 20\% (Testing) có đảm bảo tính tuần tự thời gian.

\begin{table}[h!]
\centering
\caption{Bảng So sánh Kết quả (Test Set) của các Mô hình trên 3 Tập dữ liệu}
\renewcommand{\arraystretch}{1.5}
\begin{tabular}{|l|c|c|c|c|c|c|}
\hline
\multirow{2}{*}{\textbf{Kiến trúc}} & \multicolumn{2}{c|}{\textbf{AMZN (Cổ phiếu)}} & \multicolumn{2}{c|}{\textbf{Vàng (XAU/USD)}} & \multicolumn{2}{c|}{\textbf{Bạc (XAG/USD)}} \\ \cline{2-7} 
 & \textbf{MAE} & \textbf{RMSE} & \textbf{MAE} & \textbf{RMSE} & \textbf{MAE} & \textbf{RMSE} \\ \hline
\textbf{Simple RNN} & 0.0412 & 0.0528 & 0.0385 & 0.0471 & 0.0551 & 0.0689 \\ \hline
\textbf{LSTM}       & 0.0215 & \textbf{0.0289} & \textbf{0.0198} & \textbf{0.0245} & 0.0312 & 0.0410 \\ \hline
\textbf{GRU}        & \textbf{0.0210} & 0.0295 & 0.0205 & 0.0253 & \textbf{0.0285} & \textbf{0.0375} \\ \hline
\end{tabular}
\end{table}

\subsubsection{Phân tích kết quả}
\begin{enumerate}
    \item \textbf{Amazon Stock (AMZN)}: 
    Biến động của chứng khoán Amazon phụ thuộc mạnh vào các xu hướng dài hạn. Mô hình LSTM và GRU hoàn toàn áp đảo RNN thông thường, giảm một nửa sai số. GRU có MAE tốt nhất, nhưng LSTM lại nhỉnh hơn một chút ở RMSE, cho thấy GRU có thể đã bỏ sót một vài điểm Outlier (cú sốc giá), khiến bình phương sai số bị phạt nặng hơn.
    
    \item \textbf{Vàng (Gold)}:
    Với tập dữ liệu vàng vốn dĩ có các chu kỳ dài hạn (long-term cycles) rất mạnh, LSTM thể hiện hiệu suất tốt nhất (MAE 0.0198). Cổng bộ nhớ (cell state) của LSTM giúp duy trì các mô hình giá trị dài hạn mà không bị các nhiễu nhỏ trong ngắn hạn làm "quên" đi xu hướng chính.
    
    \item \textbf{Bạc (Silver)}:
    Dữ liệu bạc có độ nhiễu cao và biến động gắt (high volatility) hơn vàng. Trong bối cảnh này, kiến trúc GRU lại vượt trội. Việc kết hợp cổng quên và cập nhật thành một giúp GRU có khả năng phản ứng (react) nhanh chóng hơn trước các biến động tức thời, cho sai số thấp nhất (RMSE 0.0375).
\end{enumerate}

\subsubsection{Kết luận về việc lựa chọn Thư viện}
\begin{itemize}
    \item Việc cài đặt \textbf{From Scratch bằng NumPy} giúp hiểu cực sâu về toán học đằng sau quá trình truyền tiến (Forward) và lan truyền ngược (BPTT). Tuy nhiên tốc độ thực thi rất chậm và dễ bị tràn số (NaN gradients).
    \item \textbf{Keras/TensorFlow} mang lại tốc độ phát triển nhanh, tính trừu tượng hóa cao, cực kỳ phù hợp cho các luồng công việc triển khai dự án công nghiệp nhanh chóng.
    \item \textbf{PyTorch} cân bằng giữa tính hiệu suất và sự mềm dẻo. Khả năng gỡ lỗi từng dòng (eager execution) và tùy chỉnh hàm chi phí hay lịch biểu học (schedulers) giúp PyTorch luôn là ưu tiên trong nghiên cứu học thuật về chuỗi thời gian.
\end{itemize}


\newpage
\section*{Tài liệu tham khảo}
\addcontentsline{toc}{section}{Tài liệu tham khảo}
\begin{enumerate}[label={[\arabic*]}]
    \item Goodfellow, I., Bengio, Y., \& Courville, A. (2016). \textit{Deep Learning}. MIT Press.
    \item Russell, S., \& Norvig, P. (2020). \textit{Artificial Intelligence: A Modern Approach}. Pearson.
    \item Các tài liệu thực hành và bài giảng môn Lập trình hệ thống thông minh.
\end{enumerate}

\end{document}
